% [TRIM3 2026-09-30] paragraphs removed from c7_discussion.tex (original text)

% Chapter intro, roadmap paragraph
This chapter summarizes each chapter, argues that no single method fits every dataset and that simple methods are often better, and closes with limitations, future directions and conclusions. % TODO: REVIEW MODIFICATION -- added the Conclusions section

% Summary, paragraph on Chapters 1 and 2
QSAR models are trained on measured labels that carry experimental noise, and most work on label noise in QSAR assumes that this noise is Gaussian. Chapter~\ref{chap:background} defined the representations, models and noise types used throughout, the NoiseInject package that injects the noise, and the procedure by which models were evaluated.

% Summary, "how the pieces fit" paragraph
Chapter~\ref{chap:robustness} tells a reader which model to choose when labels are noisy, and Chapters~\ref{chap:flaws} to~\ref{chap:similarity} tell them which features to give that model. What works depends on the dataset, as the next section argues. The contributions of this thesis are the NoiseInject package, the adaptive hybrid representation and its blueprints, and evidence that distance-based noise handling on standard representations usually made models worse. The resulting papers and software are listed in the Summary of Research Contributions.
